\scanpage{41}
\uline{Cor.3} Soit un schéma abélien sur un $S$ loc. noeth, il y a une \uline{seule} structure multiplicative possible admettant la section $e$ comme section unité bilatère.

[S'obtient aussi par d'autres voies \unsure{?} cf \unsure{\textemdash}]

À partir de l'unicité, on prouve l'existence quand on est sur un anneau \uline{local} noeth, $G$ simple et propre $S$, $\unsure{A'|a'|0'|\text{à}}$ section marquée, $G_0$ schéma abélien sur $k$ : on applique la \uline{technique} de \uline{descente}. Enfin, On doit encore passer de là au cas où $S$ est loc. noeth. connexe. On est ramené à prouver que l'ens des pts où la fibre est abélienne (qui est déjà \uline{ouvert}) est aussi \uline{fermé}. Cela ramène au cas où $S=\operatorname{Spec}(V)$ $V$ anneau de valuation discrète. Alors c'est connu. \deleted{D'où} [à développer] D'\uline{où}

\uline{Théorème} Soit $S$ loc. noeth. connexe, $G$ simple et propre sur $S$ : section marquée. Si \uline{une} fibre est abélienne il y a sur $G$ une structure mult. unique t.q. \deleted{\unsure{$e$ soit $e$|$e\times e$}}\inserted{$e$ soit section \unsure{unité}} , \deleted{et elle}\inserted{\unsure{et qu'elle|et elle}} est ``abélienne''
% archive sheet number: 40
